\documentclass[11pt]{article}
\usepackage[utf8]{inputenc}
\usepackage[T1]{fontenc}
\usepackage{lmodern}
\usepackage{amsmath,amssymb}
\usepackage{graphicx}
\usepackage{booktabs}
\usepackage{ulem}
\usepackage{fixltx2e}
\usepackage[margin=20mm,a4paper]{geometry}
\usepackage{hyperref}
\hypersetup{colorlinks=true,linkcolor=blue,urlcolor=blue}
\usepackage{kotex}
\title{My Document Converter 샘플}
\author{SHKWON}
\date{2026-09-22}

\begin{document}
\maketitle

\section*{소개 (Introduction)}\label{소개-introduction}

txt2tags 는 \textbf{굵게}, \emph{기울임}, \underline{밑줄}, \verb|코드| 와
\href{https://example.com}{링크} 를 지원합니다.
This line is in English.

\subsection*{목록 (Lists)}\label{목록-lists}

\begin{itemize}
\item 첫 번째 항목
\item 두 번째 항목
\item 세 번째 항목
\end{itemize}

\begin{enumerate}
\item 순서가 있는 항목
\item 두 번째
\item 세 번째
\end{enumerate}

\subsection*{코드 (Code)}\label{코드-code}

\begin{verbatim}
function greet(name) {
  console.log('Hello, ' + name)
}
\end{verbatim}

\subsection*{표 (Table)}\label{표-table}

\begin{table}[htbp]
\centering
\begin{tabular}{ll}
\toprule
형식 & 확장자 \\
\midrule
Markdown & .md \\
HTML & .html \\
\bottomrule
\end{tabular}
\end{table}

\end{document}
